\chapter{整除；有序域。}

初等算术运算，尤其是分数的演算，必然导致了对整数整除性的大量经验性验证。但无论是巴比伦人（尽管他们如此精于代数），还是埃及人（尽管他们对分数的计算如此灵巧），似乎都不曾知道支配这些性质的一般法则，这里的开创之功要归于希腊人。他们的算术工作——欧几里得 {[}107{]} 的第 VII 卷和第 IX 卷中有对其精湛的阐述——比起他们在其他数学分支中的那些最美的发现也毫不逊色。两个整数的最大公约数的存在性，在第 VII 卷一开头就已通过以“欧几里得算法”之名著称的程序得到证明；¹ 它是此后一切进展（素数的性质、最小公倍数的存在性与求法等等）的基础；而这一结构的顶冠则是两个出色的定理：其一是证明素数有无穷多个（第 IX 卷，命题 20），其二是给出一个从某些素数出发构造偶完全数的程序（正如欧拉后来所证明的，这个程序实际上给出了所有的偶完全数）。唯有素因子分解的存在性和唯一性没有以一般形式得到证明；不过欧几里得仍明确证明了每个整数都被某个素数整除（第 VII 卷，命题 31），以及下面两个命题（第 IX 卷，命题 13 和 14）：

“如果从一个单位开始，取任意多个数排成常数比的数列{[}即等比数列{]}，且紧接单位的那个数是素数，则最大的那个数除数列中出现的那些数外不能被任何其他数整除”（换句话说，素数的幂 \(p^n\) 只能被 \(p\) 的幂中指数 \(\leq n\) 者整除）。

“如果某数是能被{[}给定的{]}诸素数整除的数中最小的一个，那么除最初{[}给定{]}整除它的那些素数外，它不能被任何其他素数整除”（换句话说，互异素数之积 \(p_1 \ldots p_k\) 除 \(p_1, \ldots, p_k\) 外没有别的素因子）。

因此看来，欧几里得之所以没有陈述一般定理，只是由于他的术语和记法仅足以处理整数的任意幂。\(^{2}\)

尽管细致的研究使人认为，欧几里得的文本中存在着若干先后相继的层次，各自对应于算术发展的一个阶段，\(^{3}\) 但这一演化似乎在公元前 5 世纪初至公元前 4 世纪中叶之间即已完全完成，人们只能赞叹其中展现的敏锐与逻辑的稳健：要再过两千年，我们才能在算术中目睹可比拟的进展。

正是所谓“不定”问题或“丢番图”问题，处于数论此后种种进展的源头。“丢番图方程”一词按今天的用法，从历史上看并非完全有据；一般用它指整系数的代数方程（或方程组），并且只求整数解：当方程是“确定的”、也就是说只有有限多个（实数或复数）解时，这样的问题通常无解；相反，当未知数个数多于方程个数时，它往往有解。不过，尽管丢番图显然是最早考虑“不定”问题的人，他却只是例外地寻求整数解，大多数时候满足于得到一个有理数解 {[}91 a{]}。这类问题他大多能靠代数计算解决，其中未知数的算术性质根本不起作用；\(^{4}\) 因此整除理论在其中只扮演一个很不显眼的角色（“素数”一词只被提及一次（{[}91 a{]}，第 V 卷，问题 9，第 I 卷，第~334-335~页），而互素数的概念也只是在断言两个互素数的商除非其中每一个都是平方数便不可能是平方数的那条定理处才被提到）。\(^{5}\)

不定方程整数解的研究，实际上直到中世纪盛期的中国数学家和印度数学家才真正开始。前者似乎是由于编制历法中的实际问题（其中确定若干天文现象周期的公共周期，恰好构成一个一次的“丢番图”问题）而引到这类思考的；无论如何，联立线性同余方程组的一条解法规则总要归功于他们（无疑在公元 4 世纪至 7 世纪之间）。至于印度人——他们的数学从 5 世纪到 13 世纪处于全盛时期——他们不仅知道如何系统地（通过应用欧几里得算法）处理任意多个未知数的线性丢番图方程组，\(^{6}\) 而且是最早着手并解决了二次问题的人，其中包括“费马方程” \(Nx^{2} + 1 = y^{2}\) 的某些特殊情形（{[}78{]}，第 II 卷，第~87-307~页）。

我们在这里无须追述次数 \textgreater{} 1 的丢番图方程理论的历史；继费马、欧拉、拉格朗日和高斯的工作之后，这一理论在 19 世纪以代数整数理论告终（参见第~93~页以下）。正如我们已经指出的（参见第~58~页），线性方程组的研究在这一时期似乎不再提出值得注意的问题，因而多少受到忽视：特别是，人们既不去寻求任意方程组的解的一般必要条件，也不去描述解的集合。然而，大约在 19 世纪中叶，埃尔米特由于自己在数论方面的研究而用到关于线性丢番图方程的若干引理，尤其是整系数线性代换的一种“简化形”（{[}159{]}，第 I 卷，第~164~页和第~265~页）；最后，在黑格于 1858 年给出了秩等于方程个数的方程组的必要条件之后，H. J. 史密斯于 1861 年以一种“标准形”定义了整元素矩阵的不变因子（{[}287{]}，第 I 卷，第~367-409~页）。

但在此期间，阿贝尔群的概念自高斯引入（参见第~60~页）之后逐渐明确起来，这一概念在数论此后的发展中所占的分量也随之增大。在《算术研究》中对给定判别式的二次型类所成的有限阿贝尔群所作的特别深入的研究里，高斯很快察觉到这些群中有一些并不是循环群：“在这种情形”，他说，“一个基{[}也就是说一个生成元{]}是不够的，必须取两个或更多个，它们通过乘法和合成 \(^{7}\) 将能产生所有的类”（{[}124 a{]}，第 I 卷，第~374-375~页）。高斯说这些话时是否想要描述群分解为循环群的直积，这一点并不确定；不过在《算术研究》的同一页上，他证明了群中存在一个元素，其阶是所有元素的阶的最小公倍数——换句话说，他得到了该群的最大不变因子的存在性（{[}124 a{]}，第 I 卷，第~373~页）——；另一方面，直积的概念他是熟悉的，因为在一篇写于 1801 年但生前未曾发表的手稿中，他勾勒了有限阿贝尔群分解为 \(p\)-群之直积 \(^{8}\) 的一个一般证明（{[}124 a{]}，第 II 卷，第~266~页）。无论如何，1868 年，高斯著作的出版者谢林受这些结果（尤其是他刚刚重新发现的这篇手稿）的启发，证明了（仍是对二次型类群而言）一般分解定理（{[}272{]}，第 I 卷，第~135-148~页），其所用的方法两年后由克罗内克以抽象的方式重新提出（{[}186 a{]}，第 I 卷，第~273-282~页），本质上就是今天仍在使用的方法。至于无挠阿贝尔群，我们已经说过（参见第~64~页以下）椭圆函数与阿贝尔积分的理论——它由高斯、阿贝尔和雅可比发展起来——如何一步步导致认清其结构；把一个无限群分解为单演群之直和的第一个也是最著名的例子，是狄利克雷 1846 年在他关于代数数域的单位的论文中给出的（{[}92{]}，第 I 卷，第~619-644~页）。但直到 1879 年，有限型阿贝尔群的理论与史密斯定理之间的联系才由弗罗贝尼乌斯和施蒂克贝格尔明确地认识到并加以使用（{[}120{]}，§10）。

大约在同一时期，矩阵（具实系数或复系数的）相似理论也趋于完成。线性代换的特征值概念明确出现于常系数线性微分方程组的理论中，拉格朗日把它应用于小振动理论（{[}191{]}，第 I 卷，第~520~页），拉格朗日（{[}191{]}，第 VI 卷，第~655-666~页）和拉普拉斯（{[}193{]}，第 VIII 卷，第~325-366~页）把它应用于行星的“长期”不等式。它也隐含在 18 世纪中叶前后处理的许多其他问题中，比如圆锥曲线或二次曲面的轴的确定（最先由欧拉完成（{[}108 a{]}，（1），第 IX 卷，第~384~页）），又如固体的惯量主轴的研究（也由欧拉加以发展（{[}108 a{]}，（2），第 III 卷，第~200-201~页）；惯量主轴是德·塞格纳于 1755 年发现的）；今天我们知道，正是它（以隐蔽得多的形式）也参与了偏微分方程理论的发端，尤其是弦振动方程。但（且不谈后一情形）这些问题之间的共同渊源直到柯西（{[}56 a{]}，（2），第 V 卷，第~248~页和第 IX 卷，第~174~页）才勉强被认识到。此外，由于这些问题大多涉及对称矩阵，起初主要研究的是这些矩阵的特征值；这里要注意，早在 1826 年，柯西就证明了这些矩阵的特征值在相似下不变，并证明了 3 阶对称矩阵的特征值都是实数（{[}56 a{]}，（2），第 V 卷，第~248~页），三年后他把这一结果推广到任意的实对称矩阵（{[}56 a{]}，（2），第 IX 卷，第~174~页）。⁹ 默比乌斯于 1827 年引入的射影对应的一般概念（{[}223{]}，第 I 卷，第~217~页）很快引出这些变换的分类问题（首先是 2 维和 3 维情形），它不是别的，正是相应矩阵的相似问题；但在很长时间里，这个问题只用 19 世纪中叶风行的“综合”方法来处理，其进展（无论如何相当缓慢）似乎没有对特征值理论产生任何影响。几何中的另一个问题即圆锥曲线束或二次曲面束的分类则不然，从现代观点看，它归结为矩阵 \(U + \lambda V\) 的初等因子的研究，其中 \(U\) 和 \(V\) 是两个对称矩阵；西尔维斯特 1851 年正是沿着这条路线处理这个问题，他仔细考察（目的是为所考虑的束找出“标准形”）矩阵 \(U + \lambda V\) 的各子式在给 \(\lambda\) 代入一个使其行列式为零的值之后变成什么（{[}304{]}，第 I 卷，第~219-240~页）。特征值理论的纯代数方面也在同步进展；于是若干作者（其中有西尔维斯特本人）大约在 1850 年证明了 \(U^n\) 的特征值是 \(U\) 的特征值的 n 次幂；而 1858 年，凯莱在他建立矩阵演算的那篇论文中（{[}58{]}，第 II 卷，第~475-496~页），对任意阶的方阵陈述了“凯莱—哈密顿定理”，\(^{10}\) 他只满足于对 2 阶和 3 阶的矩阵通过直接计算给出证明。最后，1868 年，魏尔斯特拉斯重新采用西尔维斯特的方法，对一个“束” \(U + \lambda V\) 得到了“标准形”，这一次 \(U\) 和 \(V\) 是不必对称的方阵，唯一的限制条件是 det( \(U + \gamma V\) )

不恒为零；他由此给出任意（复元素的）方阵的初等因子的定义，并证明这些初等因子在相似的意义下刻画该方阵（{[}329 a{]}，第 II 卷，第~19-44~页）；这些结果两年后又由若尔当部分地（而且显然是独立地）重新得到 \(^{11}\)（{[}174 a{]}，第~114-125~页）。这里同样是弗罗贝尼乌斯，他于 1879 年表明，魏尔斯特拉斯定理可以容易地从推广到多项式的史密斯理论中导出（{[}119{]}，第 I 卷，第~482-544~页，§13）。

我们刚才提到了单变量多项式的整除理论；多项式除法的问题从代数的一开始就自然会提出，它是乘法运算的逆运算（后者至少对低次多项式而言已为丢番图所知）；但可以想见，在变量的各种幂的一套连贯记法确立之前，很难一般地处理这个问题。事实上，如我们所知的那种多项式“欧几里得”除法的例子，在 17 世纪中叶之前很难找到；\(^{12}\) 而 S. 斯泰芬（他本质上使用了指数记法）似乎是最先想到由此推出“欧几里得算法”的推广、用以求两个多项式的最大公因式的人（{[}295{]}，第 I 卷，第~54-56~页）。有鉴于此，直到 18 世纪中叶，整除性概念一直是有理整数所特有的。是欧拉于 1770 年开启了算术的新的一章，他不无胆量地把整除性概念推广到二次扩张中的整数：为求形如 \(x^{2}+cy^{2}\) 的数（ \(x,y,c\) 为有理整数）的因数，他写下 \(x+y\sqrt{-c}=(p+q\sqrt{-c})(r+s\sqrt{-c})\) （ \(p,q,r,s\) 为有理整数），并对两个成员取范数，从而毫不犹豫地断言，他以这种方式得到了 \(x^{2}+cy^{2}\) 的所有形如 \(p^{2}+cq^{2}\) 的因数（{[}108 a{]}，（1），第 I 卷，第~422~页）。换句话说，欧拉的论证就好像环 \(\mathbb{Z}[\sqrt{-c}]\) 是主理想环一样；稍后，他用一个类似的论证把“无穷递降”方法应用于方程 \(x^{3}+y^{3}=z^{3}\) （这归结为写下 \(p^{2}+3q^{2}\) 是一个立方数，而他通过陈述 \(p+q\sqrt{-3}=(r+s\sqrt{-3})^{3}\) 做到这一点）。但从 1773 年起，拉格朗日证明了（{[}191{]}，第 III 卷，第~695-795~页）形如 \(x^{2}+cy^{2}\) 的数的因数并非总具这种形式，这是那个基本困难的第一个例子；这一困难后来在高斯及其后继者关于单位根域中整除性的研究中，将以清晰得多的面目出现；\(^{13}\) 一般说来，不可能把有理整数的那两条本质的整除性质——最大公因数的存在和素因子分解的唯一性——直接推广到这些域。这里不是详细描述下述情形的地方：库默尔就单位根域 {[}188 b{]}、\(^{14}\) 继而戴德金和克罗内克就任意的代数数域，如何通过创立理想理论而成功地克服这一可怕的障碍——理想理论是现代代数最具决定性的进展之一（参见第~93~页以下）。但戴德金对各种数学理论的基础总是怀有好奇，并不满足于这一成功；他分析了整除关系的机制，在一篇论文中奠定了有序群的现代理论的基础，这篇论文（未曾在其同时代人中激起波澜，并被遗忘了 30 年）按年代算无疑属于关于公理化代数的最早一批工作之列（{[}79, 第 II 卷，第~103-147~页).

从 18 世纪中叶起，寻求“代数基本定理”的一个证明成了当务之急（参见第~74~页）。我们在这里无须回顾达朗贝尔的尝试，他开创了使用无穷小演算的一连串证明（参见第~161~页）。但 1749 年，欧拉以另一种方式处理这个问题（{[}108 a{]}，（1），第 VI 卷，第~78-147~页）：对每个实系数多项式 f，他试图证明分解 \(f = f_{1}f_{2}\) 的存在性，即分解为两个（非常数的）实系数多项式，这将通过对 f 的次数作归纳而给出“基本定理”的一个证明。他甚至指出，如他所言，只需在第一个奇次数的因子处停下即可，因此一切困难都归结为考虑 f 的次数 n 为偶数的情形。于是欧拉只限于研究所求的两个因子都是 n/2 次的情形，他指出，通过一个恰当进行的消元计算，\(f_{1}\) 和 \(f_{2}\) 的未知系数可以用一个实系数方程的根有理地表示出来，该方程的两端项符号相反，因而至少有一个实根。但欧拉的证明只是一个概要，其中若干本质之点被略过未论；直到 1772 年，拉格朗日才成功地解决了这个证明所提出的困难（{[}191{]}，第 III 卷，第~479-516~页），他借助一项极其冗长而极其细致的分析，在使用他（可以说是）新创造的“伽罗瓦”方法上展现出惊人的娴熟（参见第~75~页以下）。

尽管如此，拉格朗日与欧拉及其所有同时代人一样，毫不犹豫地在多项式的一个“根域”中作形式推理（也就是说，用他的语言，考虑该多项式的“虚根”）；他那个时代的数学没有为这类论证提供任何正当性。于是，对 18 世纪那种狂热的形式主义自始便自觉抱有敌意的高斯，在其博士论文中猛烈地奋起反对这种滥用（{[}124 a{]}，第 III 卷，第~3~页）。但如果他没有感觉到这里涉及的是一种表述有缺陷而其本身正确的论证，那他也就不是他自己了。于是我们看到，几年之后（{[}124 a{]}，第 III 卷，第~33~页；另见 {[}124 b{]}），他重新采用欧拉论证的一个较简单的变种——它早在 1759 年就已被德·丰瑟内提出（但后者未能把它圆满完成）——并由此推出“基本定理”的一个新证明，其中他小心避免对“虚”根的任何使用：这些根被未定元的巧妙添加和特殊化所取代。

于是，拓扑学在“基本定理”中所起的作用被归结为单独一条定理：实系数多项式不可能在一个区间内变号而不变为零（多项式情形的波尔查诺定理）。这条定理也是多项式（实系数的）实根的一切分离准则的根基，而这些准则在整个 19 世纪都属于代数最受喜爱的课题之列。\(^{15}\) 在这些研究中，人们不可能不注意到，在那里起本质作用的与其说是 R 的拓扑，不如说是 R 的序结构；\(^{16}\) 例如，多项式情形的波尔查诺定理对全部实代数数所成的域仍然成立。这一思想运动在有序域的抽象理论中达到了它的归宿，该理论由 E. 阿廷和 O. 施赖尔创立（{[}7 b{]} 与 {[}8 a and b{]}）；其更为显著的成果之一，无疑是发现了域上序关系的存在性与该域的纯粹代数性质相联系。
% semantic-fragments: legacy-input-seam
\relax{}
